\subsection{6. Ý nghĩa khoa học, thực tiễn và Bố cục luận văn}

\indent\indent \textbf{Ý nghĩa khoa học:}
\begin{itemize}
    \item Đề xuất thành công mô hình học sâu đa nhiệm \textbf{LightweightDTD} kết hợp hiệu quả giữa đặc trưng không gian và miền tần số nén JPEG.
    \item Chứng minh rằng một kiến trúc tinh gọn (~4.15 triệu tham số) hoàn toàn có thể đạt được độ chính xác phân đoạn 99.68\% và Top-3 Class Hit Rate 98.30\%, mở ra hướng tiếp cận tối ưu hóa cho các bài toán pháp y ảnh tài liệu.
\end{itemize}
\vspace{0.3cm}

\textbf{Ý nghĩa thực tiễn:}
\begin{itemize}
    \item Cung cấp công cụ tự động hóa việc xác thực tính nguyên bản của hóa đơn, hợp đồng, chứng minh nhân dân trong các dịch vụ công, ngân hàng và doanh nghiệp.
    \item Kích thước mô hình nhỏ và yêu cầu phần cứng thấp, thuận tiện tích hợp vào các thiết bị quét văn bản hoặc ứng dụng di động xử lý thời gian thực.
\end{itemize}
\vspace{0.3cm}

\textbf{Bố cục luận văn:}
Luận văn gồm 3 phần chính:
\begin{itemize}
    \item \textbf{Phần I: Phần giới thiệu:} Trình bày đặt vấn đề, các nghiên cứu liên quan, mục tiêu, đối tượng, phương pháp nghiên cứu và ý nghĩa của đề tài.
    \item \textbf{Phần II: Phần nội dung:} Gồm 3 chương:
    \begin{itemize}
        \item \textit{Chương 1: Cơ sở lý thuyết:} Trình bày nguyên lý EfficientNet-B0, biến đổi BlockDCT, mô-đun chú ý scSE và các độ đo đánh giá.
        \item \textit{Chương 2: Thiết kế và Cài đặt giải pháp:} Chi tiết kiến trúc LightweightDTD, nhánh không gian, nhánh tần số, bộ giải mã LiteMID, hàm mất mát đa nhiệm và tiền xử lý tập dữ liệu Minh-TFR.
        \item \textit{Chương 3: Thực nghiệm và Đánh giá kết quả:} Trình bày kết quả huấn luyện, các chỉ số định lượng trên tập kiểm thử và trực quan hóa kết quả dự đoán.
    \end{itemize}
    \item \textbf{Phần III: Phần kết luận và Hướng phát triển:} Tổng kết các đóng góp, phân tích hạn chế và đề xuất hướng mở rộng trong tương lai.
\end{itemize}
